\uuid{Hv7F}
\titre{Primitives usuelles}
\auteur{hanani ; mzari}
\organisation{exo7}
\licence{CC BY-NC-SA 4.0}
\datecreate{2020-01-20}
\niveau{L1}
\module{Analyse}
\chapitre{Calcul d'intégrales}
\sousChapitre{Primitives diverses}
\difficulte{3}
\theme{Primitives, primitives usuelles}
\exo7qcmid{lille-2.primitives.028}

\contenu{
\begin{question}
\qclassification{127.11}
\qtags{motcle=primitives usuelles}
Le changement de variable $u=2+\sin x$ donne :
\begin{answers}
\good{$\displaystyle \int \cos x\mathrm{e}^{2+\sin x}\mathrm{d}x=\mathrm{e}^{2+\sin x}+k$, où $k\in \Rr$.}
\bad{$\displaystyle \int \mathrm{e}^{2+\sin x}\mathrm{d}x=\int \mathrm{e}^u\mathrm{d}u$.}
\bad{$\displaystyle \int \frac{\mathrm{d}x}{2+\sin x}=\int \frac{\mathrm{d}u}u$.}
\good{$\displaystyle \int \frac{\cos x\mathrm{d}x}{2+\sin x}=\ln(2+\sin x)+k$, où $k\in \Rr$.}
\end{answers}
\begin{explanations}
Avec $u=2+\sin x$, on a : $\mathrm{d}u=\cos x\mathrm{d}x$ et puis
$$\int \cos x\mathrm{e}^{2+\sin x}\mathrm{d}x=\int \mathrm{e}^{u}\mathrm{d}u=\mathrm{e}^{u}+k=\mathrm{e}^{2+\sin x}+k,\; k\in \Rr.$$
Ensuite, $\displaystyle \int \frac{\cos x\mathrm{d}x}{2+\sin x}=\int \frac{\mathrm{d}u}{u}=\ln |u|+k=\ln (2+\sin x)+k$, $k\in \Rr$.
\end{explanations}
\end{question}
}
